% preamble-exam-zh.tex — 共享 preamble
% 用于所有 exam-zh 模板
%
% 样式策略：
%   屏幕 → 语义彩色（蓝=关键想法，红=易错，绿=路线，灰=提示/教师备注）
%   打印 → 自动降级为黑白（通过 print mode 或 monochrome 选项）

\documentclass{exam-zh}

% ---------- 基础配置 ----------
\examsetup{
  page/size = a4paper,
  page/show-head = false,
  page/show-foot = false,
  sealline/show = false,
  font = times,
  math-font = xits,
}

\usepackage{tcolorbox}
\usepackage{needspace}
\tcbuselibrary{skins,breakable}

% ---------- 打印/屏幕颜色切换 ----------
% 用法：\documentclass[monochrome]{exam-zh} 即可切换为纯黑白
% 注意：不要使用 \if@xxx 放在 makeatother 外部；这里改成普通条件名。
\newif\ifeduprintcolor
\eduprintcolortrue

\makeatletter
\@ifclasswith{exam-zh}{monochrome}{\eduprintcolorfalse}{}
\makeatother

% ---------- 语义颜色定义 ----------
% 屏幕：有色彩区分；打印：降级为灰度
\ifeduprintcolor
  % --- 屏幕彩色模式 ---
  \definecolor{edu-blue}{RGB}{47, 82, 122}       % 关键想法
  \definecolor{edu-blue-bg}{RGB}{243, 246, 252}
  \definecolor{edu-red}{RGB}{180, 40, 40}         % 易错提醒
  \definecolor{edu-red-bg}{RGB}{255, 245, 240}
  \definecolor{edu-green}{RGB}{34, 100, 60}       % 路线图
  \definecolor{edu-green-bg}{RGB}{242, 250, 245}
  \definecolor{edu-orange}{RGB}{160, 90, 0}       % 训练目标
  \definecolor{edu-orange-bg}{RGB}{255, 248, 235}
  \definecolor{edu-gray}{RGB}{100, 100, 100}      % 提示/教师备注
  \definecolor{edu-gray-bg}{RGB}{248, 248, 248}
  \definecolor{edu-purple}{RGB}{90, 50, 120}      % 思考题
  \definecolor{edu-purple-bg}{RGB}{248, 244, 252}
\else
  % --- 打印黑白模式 ---
  \definecolor{edu-blue}{gray}{0.15}
  \definecolor{edu-blue-bg}{gray}{0.96}
  \definecolor{edu-red}{gray}{0.25}
  \definecolor{edu-red-bg}{gray}{0.97}
  \definecolor{edu-green}{gray}{0.20}
  \definecolor{edu-green-bg}{gray}{0.97}
  \definecolor{edu-orange}{gray}{0.25}
  \definecolor{edu-orange-bg}{gray}{0.98}
  \definecolor{edu-gray}{gray}{0.40}
  \definecolor{edu-gray-bg}{gray}{0.97}
  \definecolor{edu-purple}{gray}{0.30}
  \definecolor{edu-purple-bg}{gray}{0.98}
\fi

% ---------- 自定义教学环境 ----------
% 共性：enhanced + breakable（跨页不断裂）

% 原题卡片 — 强边框，突出显示
\newtcolorbox{problemcard}{
  enhanced, breakable,
  colback=edu-blue-bg,
  colframe=edu-blue,
  boxrule=1.2pt,
  arc=0pt,
  left=3mm, right=3mm, top=3mm, bottom=3mm,
}

% 解题路线图 — 左边线 + 浅底
\newtcolorbox{routemap}{
  enhanced, breakable,
  colback=edu-green-bg,
  colframe=edu-green!30,
  boxrule=0pt,
  borderline west={2.5pt}{0pt}{edu-green},
  left=3mm, right=3mm, top=3mm, bottom=3mm,
}

% 关键想法 — 蓝色左边线
\newtcolorbox{keyideabox}{
  enhanced, breakable,
  colback=edu-blue-bg,
  colframe=edu-blue!30,
  boxrule=0pt,
  borderline west={2.5pt}{0pt}{edu-blue},
  left=3mm, right=3mm, top=3mm, bottom=3mm,
}

% 易错提醒 — 红色左边线
\newtcolorbox{mistakebox}{
  enhanced, breakable,
  colback=edu-red-bg,
  colframe=edu-red!20,
  boxrule=0pt,
  borderline west={3pt}{0pt}{edu-red},
  left=3mm, right=3mm, top=2.5mm, bottom=2.5mm,
}

% 提示 — 灰色虚线左边线
\newtcolorbox{hintbox}{
  enhanced, breakable,
  colback=edu-gray-bg,
  colframe=edu-gray!20,
  boxrule=0pt,
  borderline west={2pt}{0pt}{edu-gray, dashed},
  left=3mm, right=3mm, top=2mm, bottom=2mm,
}

% 思考题 — 紫色虚线左边线
\newtcolorbox{thinkbox}{
  enhanced, breakable,
  colback=edu-purple-bg,
  colframe=edu-purple!20,
  boxrule=0pt,
  borderline west={2pt}{0pt}{edu-purple, dashed},
  left=2.5mm, right=2.5mm, top=2.5mm, bottom=2.5mm,
}

% 教师备注 — 灰色左边线，小字
\newtcolorbox{teachernote}{
  enhanced, breakable,
  colback=edu-gray-bg,
  colframe=edu-gray!15,
  boxrule=0pt,
  borderline west={2pt}{0pt}{edu-gray!60},
  left=2.5mm, right=2.5mm, top=2.5mm, bottom=2.5mm,
  fontupper=\small,
  colupper=black!60,
}

% 训练目标 — 橙色左边线，小字
\newtcolorbox{traininggoal}{
  enhanced, breakable,
  colback=edu-orange-bg,
  colframe=edu-orange!15,
  boxrule=0pt,
  borderline west={1.5pt}{0pt}{edu-orange},
  left=2mm, right=2mm, top=1mm, bottom=1mm,
  fontupper=\small,
}

% 答题区 — 无边框，纯留白
\newcommand{\answerarea}[1][40mm]{%
  \vspace{2mm}%
  \begin{tcolorbox}[
    enhanced, breakable,
    colback=white,
    colframe=white,
    boxrule=0pt,
    height=#1,
    valign=top,
    bottom=0mm,
    after skip=0mm,
  ]
  \end{tcolorbox}%
}

% ---------- 字体设置 ----------
% exam-zh (ctexbook) already sets CJK fonts via ctex.
% Only override if specific fonts are needed.
% macOS: Songti SC, STHeiti, STFangsong
% Linux: SimSun, SimHei, FangSong

% ---------- 页面设置 ----------
% exam-zh (ctexbook) already loads geometry.
% Override margins if needed:
% \geometry{top=16mm, bottom=16mm, left=14mm, right=14mm}

\examsetup{
  question/show-answer = true,
  fillin/show-answer = true,
  solution/show-solution = show-stay,
}

% ---------- 教师版解答样式 ----------
\newtcolorbox{solutionblock}{
  enhanced, breakable,
  colback=edu-blue-bg,
  colframe=edu-blue!25,
  boxrule=0pt,
  borderline west={2pt}{0pt}{edu-blue!50},
  arc=0pt,
  left=3mm, right=3mm, top=2mm, bottom=2mm,
  before skip=2mm,
  after skip=3mm,
  fontupper=\small,
}

% 解答步骤编号
\newcounter{solstep}
\newcommand{\solstep}[2]{%
  \stepcounter{solstep}%
  \par\medskip
  \noindent\hangindent=6mm
  {\color{edu-blue}\bfseries\textcircled{\scriptsize\thesolstep}\enspace #1}\enspace #2\par
}

\begin{document}

\section*{垂径定理专题练习（教师版）}

\section*{一、填空题（每题 4 分，共 20 分）}


\begin{question}[points=4]
  在 $\odot O$ 中，弦 $AB = 8$，圆心 $O$ 到 $AB$ 的距离为 $3$，则 $\odot O$ 的半径为\fillin。
  \fillin
  \paren[5]
  \begin{solutionblock}
    作 $OM \perp AB$ 于 $M$，则 $AM = 4$，在 $\mathrm{Rt}\triangle OAM$ 中 $r = \sqrt{3^2+4^2} = 5$。
  \end{solutionblock}
\end{question}



\begin{question}[points=4]
  拱桥的桥拱是圆弧形，跨度（弦长）$AB = 12$ m，拱高 $CD = 2$ m，则桥拱所在圆的半径为\fillin m。
  \fillin
  \paren[10]
  \begin{solutionblock}
    $AD = 6$，$OD = r - 2$。$r^2 = 6^2 + (r-2)^2 = 36 + r^2 - 4r + 4$，$4r = 40$，$r = 10$。
  \end{solutionblock}
\end{question}



\begin{question}[points=4]
  在 $\odot O$ 中，$AB \parallel CD$，$AB = 6$，$CD = 8$，半径 $r = 5$，则 $AB$ 与 $CD$ 的距离为\fillin。
  \fillin
  \paren[1 或 7]
  \begin{solutionblock}
    $OM = \sqrt{5^2-3^2} = 4$，$ON = \sqrt{5^2-4^2} = 3$。同侧 $= 4-3 = 1$，异侧 $= 4+3 = 7$。
  \end{solutionblock}
\end{question}



\begin{question}[points=4]
  $\odot O$ 的半径为 $10$，弦 $AB = 16$，点 $P$ 在弦 $AB$ 上且 $AP = 3$，则 $OP$ 的长为\fillin。
  \fillin
  \paren[$\sqrt{61}$]
  \begin{solutionblock}
    作 $OM \perp AB$，$AM = 8$，$OM = \sqrt{10^2-8^2} = 6$。$P$ 在 $AB$ 上，$AP = 3$，$PM = |8-3| = 5$。$OP = \sqrt{6^2+5^2} = \sqrt{61}$。
  \end{solutionblock}
\end{question}



\begin{question}[points=4]
  如图，在 $\odot O$ 中，直径 $AB$ 垂直于弦 $CD$ 于点 $H$，$AH = 2$，$CD = 8$，则 $\odot O$ 的半径为\fillin。
  \fillin
  \paren[5]
  \begin{solutionblock}
    直径 $\perp$ 弦，$DH = 4$。设 $r$，$OH = r - 2$。$r^2 = (r-2)^2 + 4^2$，$4r = 20$，$r = 5$。
  \end{solutionblock}
\end{question}


\section*{二、解答题（每题 10 分，共 30 分）}

\needspace{8\baselineskip}

\begin{problem}[points=10]
  如图，$\odot O$ 的直径 $AB \perp$ 弦 $CD$ 于 $H$，$AH = 3$，$DH = 4$，$E$ 是 $\widehat{AC}$ 上一点，联结 $AE$、$DE$。\\
(1) 求 $\odot O$ 的半径；\\
(2) 若 $AE = 5$，求 $DE$ 的长。

  \answerarea[80mm]
  \setcounter{solstep}{0}
  \begin{solutionblock}
    \textbf{答案：}(1) $r = \frac{25}{6}$；(2) 需联结 $CE$ 和 $AD$，利用 $\triangle ADE \sim \triangle ACB$，$DE = \frac{20}{3}$。\par\medskip
    \solstep{第(1)问}{$r^2 = (r-3)^2 + 4^2$，$6r = 25$，$r = 25/6$。}
    \solstep{第(2)问}{联结 $AD$。$AD = \sqrt{3^2+4^2} = 5$。$\angle ADE = \angle ABD$，$\triangle ADE \sim \triangle ABD$。$\frac{DE}{AD} = \frac{AE}{AB}$。$AB = 2r = 25/3$。$\frac{DE}{5} = \frac{AE}{25/3}$。$DE = AE \times 15/25 = 3AE/5$。}
  \end{solutionblock}
\end{problem}


\needspace{8\baselineskip}

\begin{problem}[points=10]
  已知 $AB$ 是 $\odot O$ 的直径，弦 $CD$ 交 $AB$ 于点 $E$，$\angle CEA = 45^\circ$，$OE = 4$，$DE = 3 + 2\sqrt{2}$，求弦 $CD$ 的长及 $\odot O$ 的半径。

  \answerarea[80mm]
  \setcounter{solstep}{0}
  \begin{solutionblock}
    \textbf{答案：}$CD = 6$，半径 $= \sqrt{17}$\par\medskip
    \solstep{作辅助线}{过 $O$ 作 $OM \perp CD$ 于 $M$，联结 $OD$。}
    \solstep{45°特殊角求OM和EM}{$\angle OEM = 45^\circ$，$OE = 4$。等腰直角三角形：$OM = EM = 2\sqrt{2}$。}
    \solstep{求DM和CD}{$DM = DE - EM = (3+2\sqrt{2}) - 2\sqrt{2} = 3$。$CD = 2 \times 3 = 6$。}
    \solstep{求半径}{$OD = \sqrt{(2\sqrt{2})^2 + 3^2} = \sqrt{8+9} = \sqrt{17}$。}
  \end{solutionblock}
\end{problem}


\needspace{8\baselineskip}

\begin{problem}[points=10]
  已知 $\odot O$ 的半径为 $10$，弦 $AB = 12$，弦 $CD = 16$，$AB \parallel CD$。\\
(1) 求 $AB$ 与 $CD$ 之间的距离；\\
(2) 若在 $AB$ 与 $CD$ 之间画一条与它们平行的弦 $EF$，使得 $EF$ 到 $AB$ 和 $CD$ 的距离相等，求 $EF$ 的长。

  \answerarea[80mm]
  \setcounter{solstep}{0}
  \begin{solutionblock}
    \textbf{答案：}(1) $2$ 或 $14$；(2) 同侧 $EF = 2\sqrt{51}$，异侧 $EF = 6\sqrt{11}$。\par\medskip
    \solstep{分别求弦心距}{$OM = \sqrt{100-36} = 8$，$ON = \sqrt{100-64} = 6$。}
    \solstep{分类讨论距离}{同侧 $= 8-6 = 2$，异侧 $= 8+6 = 14$。}
    \solstep{求EF弦心距}{同侧：$d_{EF} = (8+6)/2 = 7$；异侧：$d_{EF} = (8-6)/2 = 1$。}
    \solstep{求EF弦长}{同侧：$EF = 2\sqrt{100-49} = 2\sqrt{51}$；异侧：$EF = 2\sqrt{100-1} = 6\sqrt{11}$。}
  \end{solutionblock}
\end{problem}


\section*{★ 提高题（不计入总分）}

\needspace{8\baselineskip}

\begin{problem}[points=0]
  在 $\odot O$ 中，弦 $AB = 10$，弦 $CD = 10$，$AB \parallel CD$，半径 $r = \sqrt{34}$。$P$ 是 $\odot O$ 上一点（不在 $AB$、$CD$ 上），求 $\triangle PAB$ 与 $\triangle PCD$ 面积之和的最大值。

  \answerarea[80mm]
  \setcounter{solstep}{0}
  \begin{solutionblock}
    \textbf{答案：}$10\sqrt{34}$\par\medskip
    \solstep{求弦心距和两弦位置}{$d = \sqrt{34-25} = 3$。两弦等长不同弦，必在异侧，距离 $= 6$。}
    \solstep{表达面积之和}{$S = \frac{1}{2} \times 10 \times h_1 + \frac{1}{2} \times 10 \times h_2 = 5(h_1+h_2)$。}
    \solstep{求h1+h2最大值}{$h_1 + h_2$ 最大 = 直径 $= 2\sqrt{34}$。最大 $S = 10\sqrt{34}$。}
  \end{solutionblock}
\end{problem}



\end{document}
